\documentclass[tikz,border=8pt]{standalone}

\usepackage{amsmath}
\usepackage{tikz}
\usetikzlibrary{arrows.meta,calc,decorations.pathreplacing}

\begin{document}

\begin{tikzpicture}[
    >=Stealth,
    line width=0.8pt,
    every node/.style={font=\small},
]

% ============================================================
% (a) POINT CHARGE
% ============================================================

\begin{scope}[xshift=0cm,yshift=6.2cm]

    % Charge q
    \fill (0,0) circle (1.3pt);
    \node[below left] at (0,0) {$q$};

    % Point P
    \fill (1.65,2.35) circle (1.3pt);
    \node[above] at (1.65,2.35) {$P$};

    % r_QP
    \draw[->] (0.04,0.05) -- (1.65,2.35)
        node[midway,left=2pt] {$\mathbf r_{qP}$};

    % Caption
    \node[below] at (0.8,-0.65) {Point charge};
    \node[below] at (0.8,-1.00) {(a)};

\end{scope}


% ============================================================
% (b) LINE CHARGE
% ============================================================

\begin{scope}[xshift=6.4cm,yshift=6.2cm]

    % Line charge
    \draw (0.0,0.0) -- (0.75,3.15);

    % L labels
    \node[below left] at (0.0,0.0) {$L$};
    \node[above right] at (0.75,3.15) {$\lambda$};

    % Positive charges along line
    \foreach \y in {0.35,0.68,1.01,1.34,1.67,2.00,2.33,2.66} {
        \node at ({0.75*\y/3.15 - 0.08},\y) {$+$};
    }

    % Differential element
    \coordinate (Q) at (0.39,1.60);

    % dq bracket/marker
    \draw[thick] ($(Q)+(-0.06,-0.12)$) -- ($(Q)+(0.01,0.17)$);
	\node[above left] at (0.34,1.55) {$dl$}

    % P
    \fill (2.55,2.55) circle (1.3pt);
    \node[above right] at (2.55,2.55) {$P$};

    % r_QP
    \draw[->] ($(Q)+(0.02,0.02)$) -- (2.55,2.55)
        node[midway,above=1pt] {$\mathbf r_{qP}$};

    % Differential charge equations
    \node[align=left] at (1.75,1.00) {
        $dq=\lambda\,dl$\\[-1pt]
        $\ \ q=\displaystyle\int_L\lambda\,dl$
    };

    % Caption
    \node[below] at (1.25,-0.65) {Line charge};
    \node[below] at (1.25,-1.00) {(b)};

\end{scope}


% ============================================================
% (c) SURFACE CHARGE
% ============================================================

\begin{scope}[xshift=0cm,yshift=0cm]

\coordinate (A) at (-0.25,0.25);
\coordinate (B) at (1.65,1.25);
\coordinate (C) at (1.65,3.25);
\coordinate (D) at (-0.25,2.25);

\draw (A) -- (B) -- (C) -- (D) -- cycle;
% Surface label
\node at (0.25,0.70) {$S$};

% Surface charge density
\node[above left] at (-0.25,2.25) {$\sigma$};

% 4x4 grid of charges aligned with the parallelogram
% The parallelogram has basis vectors:
% v1 = B - A = (1.90, 1.00)
% v2 = D - A = (0.00, 2.00)
% So any point = A + u*v1 + v*v2, where u,v in [0,1]
\foreach \i in {1,2,3,4} {
    \foreach \j in {1,2,3,4} {
        % u and v go from 0.125 to 0.875 (centered in each cell)
        \pgfmathsetmacro{\u}{(\i-0.5)/4}
        \pgfmathsetmacro{\v}{(\j-0.5)/4}
        
        % Compute position: P = A + u*(B-A) + v*(D-A)
        \pgfmathsetmacro{\px}{-0.25 + \u*1.90 + \v*0.00}
        \pgfmathsetmacro{\py}{0.25 + \u*1.00 + \v*2.00}
        
        \node at (\px,\py) {$+$};
    }
}

% Differential surface element dS (small parallelogram aligned with S)
\coordinate (dS_center) at (0.65,1.75);

% Scale down the parallelogram by factor s
\pgfmathsetmacro{\s}{0.18}

% The parallelogram corners relative to center:
% A: (-0.25, 0.25) -> relative to center: (-0.90, -1.50)
% B: (1.65, 1.25)  -> relative: (1.00, -0.50)
% C: (1.65, 3.25)  -> relative: (1.00, 1.50)
% D: (-0.25, 2.25) -> relative: (-0.90, 0.50)

\coordinate (dA) at ($(dS_center) + \s*(-0.90, -1.50)$);
\coordinate (dB) at ($(dS_center) + \s*(1.00, -0.50)$);
\coordinate (dC) at ($(dS_center) + \s*(1.00, 1.50)$);
\coordinate (dD) at ($(dS_center) + \s*(-0.90, 0.50)$);

% Draw dS as a small parallelogram
\draw[thick] (dA) -- (dB) -- (dC) -- (dD) -- cycle;


% dS label
\node[above right] at ($(dC)+(0.05,0.05)$) {$dS$};

% P
\fill (3.10,2.55) circle (1.3pt);
\node[above right] at (3.10,2.55) {$P$};

% r_QP
\draw[->] ($(dS_center)+(0.0,0.0)$) -- (3.10,2.55)
    node[above right] at ($(dS_center)!0.5!(3.10,2.55)$) {$\mathbf r_{qP}$};

% Equations
\node[align=left] at (0.00,3.55) {
    $dq=\sigma\,dS$\\[-1pt]
    $q=\displaystyle\int_S\sigma\,dS$
};

% Caption
\node[below] at (0.70,-0.35) {Surface charge};
\node[below] at (0.70,-0.70) {(c)};

\end{scope}


% ============================================================
% (d) VOLUME CHARGE
% ============================================================

\begin{scope}[xshift=6.4cm,yshift=0cm]

    % Bean-shaped volume boundary
    \draw
        (1.00,0.40)
        .. controls (0.20,0.30) and (-0.10,0.90) .. (0.15,1.40)
        .. controls (0.30,1.70) and (0.60,2.00) .. (1.00,2.20)
        .. controls (1.40,2.40) and (2.00,2.55) .. (2.40,2.35)
        .. controls (2.80,2.15) and (2.90,1.70) .. (2.80,1.30)
        .. controls (2.70,0.90) and (2.40,0.60) .. (2.00,0.48)
        .. controls (1.60,0.38) and (1.20,0.42) .. cycle;

    % V label
    \node at (2.15,0.75) {$V$};

    % Additional charges to fill the bean shape
    \foreach \p in {
		(0.75,1.50),
        (0.65,1.80),
        (0.90,2.00),
		(2.0,2.10),
        (2.25,1.90),
	   (1.9,0.85),
	   (0.70,0.65),
        (0.50,0.80)
    } {
        \node at \p {$+$};
    }

    % Differential volume element
    \coordinate (dV) at (1.20,1.50);

 % Small cube indication of dV (3D perspective)
\draw ($(dV)+(-0.15,-0.15)$) rectangle ($(dV)+(0.15,0.15)$);
\draw ($(dV)+(-0.15,-0.15)+(0.08,0.08)$) rectangle ($(dV)+(0.15,0.15)+(0.08,0.08)$);
\draw ($(dV)+(-0.15,-0.15)$) -- ($(dV)+(-0.15,-0.15)+(0.08,0.08)$);
\draw ($(dV)+(0.15,-0.15)$) -- ($(dV)+(0.15,-0.15)+(0.08,0.08)$);
\draw ($(dV)+(0.15,0.15)$) -- ($(dV)+(0.15,0.15)+(0.08,0.08)$);
\draw ($(dV)+(-0.15,0.15)$) -- ($(dV)+(-0.15,0.15)+(0.08,0.08)$);
\node[right] at ($(dV)+(0.18,-0.05)$) {$dV$};
    % P
    \fill (3.00,2.65) circle (1.3pt);
    \node[above right] at (3.00,2.65) {$P$};

    % r_QP
    \draw[->] ($(dV)+(0.0,0.0)$) -- (3.00,2.65)
        node[midway,above] {$\mathbf r_{qP}$};

    % Equations
    \node[align=left] at (0.15,2.80) {
        $q=\displaystyle\int_V\rho\,dV$\\[-1pt]
        $dq=\rho\,dV$
    };

    % Caption
    \node[below] at (1.25,0.00) {Volume charge};
    \node[below] at (1.25,-0.35) {(d)};

\end{scope}

\end{tikzpicture}

\end{document}
